\section{Bối cảnh nghiên cứu}
\label{sec:background}

Thực khuẩn thể (bacteriophage), thường được gọi tắt là phage, là các virus xâm nhiễm và nhân lên bên trong tế bào vi khuẩn. Với số lượng ước tính lên tới $10^{31}$ cá thể~\cite{suttle2007marine}, thực khuẩn thể là dạng sống phong phú nhất trên Trái Đất và đóng vai trò trung tâm trong việc điều hòa quần thể vi khuẩn, duy trì cân bằng sinh thái vi sinh vật, và thúc đẩy quá trình chuyển gen ngang (horizontal gene transfer) giữa các vi khuẩn~\cite{koskella2014bacteria}.

Thực khuẩn thể được phân loại thành hai nhóm chính dựa trên chiến lược nhân lên~\cite{wu2021deephage}. Thực khuẩn thể độc lực (virulent phage), còn gọi là phage phân giải nghiêm ngặt (strictly lytic phage), chỉ thực hiện chu trình phân giải: sau khi xâm nhập tế bào vi khuẩn, phage chiếm đoạt bộ máy tổng hợp của tế bào chủ để sản xuất các virion con, rồi gây phá vỡ tế bào để giải phóng thế hệ phage mới~\cite{mcnair2012phacts}. Ngược lại, thực khuẩn thể ôn hòa (temperate phage) có thể thực hiện cả chu trình phân giải lẫn chu trình lysogen: trong chu trình lysogen, vật liệu di truyền của phage được tích hợp vào nhiễm sắc thể vi khuẩn chủ dưới dạng prophage, tồn tại ở trạng thái không hoạt động và nhân lên cùng hệ gen vi khuẩn qua nhiều thế hệ~\cite{howard2017lysogeny}. Prophage có thể được kích hoạt bởi các yếu tố môi trường như tia UV hoặc tác nhân gây tổn thương DNA, chuyển sang chu trình phân giải~\cite{feiner2015new}.

Việc phân loại chính xác lối sống của thực khuẩn thể có ý nghĩa quan trọng trong cả nghiên cứu vi sinh vật học và ứng dụng y sinh học. Trong nghiên cứu hệ vi sinh vật, xác định lối sống phage giúp làm sáng tỏ các tương tác phức tạp giữa phage và vi khuẩn chủ, từ đó hiểu rõ hơn vai trò sinh thái của phage trong các cộng đồng vi sinh vật~\cite{shang2023phatyp}. Trong bối cảnh kháng kháng sinh ngày càng gia tăng, liệu pháp phage (phage therapy) đang được xem xét như một giải pháp thay thế đầy hứa hẹn cho kháng sinh truyền thống~\cite{gorski2016phage,azimi2019phage}. Tuy nhiên, chỉ có thực khuẩn thể độc lực mới phù hợp cho liệu pháp phage, vì thực khuẩn thể ôn hòa có thể tích hợp vào hệ gen vi khuẩn bệnh nguyên và tiềm ẩn nguy cơ truyền các gen độc lực hoặc gen kháng kháng sinh thông qua quá trình lysogenic conversion~\cite{cobian2016viruses}.

Các phương pháp truyền thống để xác định lối sống phage dựa vào nuôi cấy trong phòng thí nghiệm và quan sát kiểu hình, nhưng tốn thời gian, đòi hỏi nhiều công sức, và không thể áp dụng cho phần lớn các phage chưa được nuôi cấy được trong các mẫu môi trường~\cite{shang2023phatyp}. Những hạn chế này đã thúc đẩy sự phát triển của các phương pháp tính toán để dự đoán lối sống phage trực tiếp từ dữ liệu trình tự gen. Với sự phát triển của công nghệ giải trình tự thế hệ mới, metagenomics đã trở thành công cụ quan trọng để nghiên cứu đa dạng sinh học của phage mà không cần nuôi cấy~\cite{mokili2012metagenomics,hayes2017metagenomic}. Trong phân tích metagenomic, hệ gen phage thường được phục hồi dưới dạng các contig—các chuỗi liên tục được lắp ráp từ các đoạn đọc ngắn chồng lấp~\cite{setubal2021metagenome,ghurye2016metagenomic}. Các contig này thường có độ dài biến thiên lớn và bị phân mảnh do độ phủ giải trình tự không đồng đều, các hiện vật kỹ thuật, và các vùng gen lặp lại, khiến việc phân loại chính xác các chuỗi không hoàn chỉnh trở thành một thách thức tính toán quan trọng.
